\section{Related Work}
\label{sec:related}

\paragraph{Deep learning on the 12-lead ECG.} PTB-XL~\cite{wagner2020ptbxl}, released through PhysioNet~\cite{goldberger2000physionet}, made a large, cardiologist-annotated 12-lead dataset with fixed patient-level folds publicly available. Strodthoff et al.~\cite{strodthoff2021benchmark} benchmarked convolutional, recurrent and feature-based models on it; for the five diagnostic superclasses, the best deep models reached macro AUROC between 0.921 and 0.930, and a wavelet-feature network 0.874. ResNet-style one-dimensional convolutional networks~\cite{he2016resnet} performed consistently well, and recurrent models were competitive. Ribeiro et al.~\cite{ribeiro2020automatic} trained a residual network on more than two million 12-lead ECGs and reported performance comparable to or better than that of cardiology residents for six abnormalities. We use the PTB-XL superclass task and folds so that our 12-lead results can be compared with this benchmark.

\paragraph{Reduced-lead and single-lead ECG analysis.} Hannun et al.~\cite{hannun2019cardiologist} showed that a deep network can classify twelve rhythm classes from single-lead ambulatory ECGs. The PhysioNet/Computing in Cardiology Challenge 2021, ``Will Two Do?''~\cite{reyna2021willtwodo}, asked explicitly how reduced lead sets (twelve, six, four, three and two leads) compare for multi-label diagnosis across several databases, including PTB-XL and the Chapman-Shaoxing/Ningbo data~\cite{zheng2020chapman}. Closest to our single-lead question, Davarmanesh et al.~\cite{davarmanesh2024leadI} detected MI from lead I of PTB-XL and reported a test AUC of 0.92 for MI versus normal records. These studies report how well reduced leads classify; they do not ask, with a pre-specified criterion, whether a reduced lead set is good enough to \emph{replace} the 12-lead ECG for each diagnosis, which is our focus.

\paragraph{Transfer and self-supervised learning.} Transfer of representations from a data-rich setting to a data-poor one is standard in deep learning; for ECGs, Mehari and Strodthoff~\cite{mehari2022ssl} showed that self-supervised pre-training on 12-lead data improves downstream PTB-XL performance and robustness. Masked modelling, introduced for language in BERT~\cite{devlin2019bert} and for images in masked autoencoders~\cite{he2022mae}, is a natural pretext task for Transformers~\cite{vaswani2017attention,dosovitskiy2021vit}. We test whether a simple masked-signal pre-training of an ECG Vision Transformer helps when only 10\% of the labels are available, and whether supervised transfer from 12 leads helps reduced-lead models.

\paragraph{Calibration, statistics and interpretability.} Modern networks are often miscalibrated, and temperature scaling fitted on a validation set is a simple remedy~\cite{guo2017calibration}. We report calibration alongside discrimination, use bootstrap confidence intervals~\cite{efron1993bootstrap} with Holm correction for multiple comparisons~\cite{holm1979}, and use Grad-CAM~\cite{selvaraju2017gradcam} and attention maps only as illustrative, not causal, explanations.
